\section{Target and horizon}
\label{sec:target}

\subsection{Timing}

Everything is computed at the \textbf{close of day $t$} from information public by then. The gate updates
every trading day, every stock in the universe gets a forecast every day, and the target is the return
from the close of $t$ to the close of $t+5$.

\begin{whybox}
Daily closing data is what CRSP provides, and using one timestamp for features, gate and target makes
look-ahead easy to rule out: every input at $t$ is computed from data up to the close of $t$, and the
target starts after it. The convention assumes trading at the close at which the signal is computed,
which is slightly optimistic; signals concentrated overnight should be checked with a one-day execution
lag.
\end{whybox}

\subsection{Horizon: five trading days}

\dec{2026-10-05, Q21, option B} Daily gate, daily cross-section, \textbf{$h=5$}. Consecutive targets overlap by
four days. The options compared were (A) $h=1$, (B) $h=5$, (C) a monthly target with a daily gate,
(D) monthly gate and cross-section.

\begin{whybox}
\begin{enumerate}
\item \textbf{Enough independent observations.} A monthly horizon leaves too few independent periods to train
  regime-specific experts: in 2015--2024 about 118, of which about 59 in high volatility, and a 5\% rank IC
  at $h=21$ was not significant. Weekly gives about five times as many.
\item \textbf{The regime effect is clearest at this horizon.} A descriptive check (2015--2024, signal
  diagnostics, not model results): one-week reversal had a rank IC of 3.26\% in high-volatility periods
  against 1.29\% in low ones. Short-term reversal is pay for providing liquidity, and it rises when markets
  are turbulent (Nagel 2012). So at a weekly horizon the cross-section demonstrably depends on the market
  state, which is precisely what a regime gate is meant to exploit.
\item \textbf{Less noise than one day.} A one-day return is dominated by noise and bid--ask bounce (the
  closing price jumping between bid and ask, which creates fake reversal; Roll 1984). One-day
  machine-learning strategies on the S\&P 500 have largely decayed since 2010 (Krauss, Do \& Huck 2017;
  Fischer \& Krauss 2018). Five days averages much of this out while keeping short-horizon effects.
\item \textbf{It respects the regime timescale.} Daily regimes typically last 20--50 days, so a 5-day window
  rarely contains a switch, and ``one regime per target window'' is a reasonable approximation. A monthly
  window is as long as a stressed regime and would break it.
\item \textbf{The gate is best identified on daily data}: thousands of daily observations instead of a few
  hundred monthly ones, and daily volatility clustering is what separates the regimes.
\item \textbf{It was already built and audited}: overlap-corrected statistics, purge and a horizon-consistent
  portfolio.
\end{enumerate}
\giveup{direct comparability with the monthly benchmark of Gu, Kelly \& Xiu; independence of
consecutive targets (statistics are overlap-corrected, and the training likelihood is used for
estimation, not inference); and lower turnover than a monthly strategy.}
\end{whybox}

\paragraph{Other horizons (discussed 8 October).} A \textbf{one-day} horizon has real advantages (no overlap,
an exact one-step gate weight, the freshest signals, closer to industry practice) but most of its edge
disappears with a realistic one-day execution delay and five times the turnover. A \textbf{three-day}
horizon is a middle ground, but adds noise and turnover, and its windows sometimes contain a weekend
and sometimes not; there is no evidence it beats five. \dec{2026-10-08, Q21 follow-up} $h=5$ stays primary,
with two robustness additions: an \textbf{IC decay curve} (rank IC of the main signals at
$h=1,2,3,5,10$, on 2000--2006 only) and a \textbf{horizon profile} ($h=1,3,5,10$) for the final model.

\begin{whybox}[Why add these rather than switch]
Whether a shorter horizon is better depends on how fast the signals' predictive power fades. That is a
measurable fact, so it should be measured, not argued; and measured on years before the pilot and the test
period, so that the horizon cannot be chosen on test results. The horizon profile for the final model shows
whether conclusions depend on the horizon, without multiplying the main study.
\end{whybox}

\subsection{Target: the market-neutral return}

\dec{2026-10-01, Q25} For stock $i$ in the universe $\mathcal S_t$, with forward log return
$r_{i,t\to t+5}=\sum_{d=1}^{5}\log(1+r_{i,t+d})$,
\begin{equation}
y_{i,t} \;=\; r_{i,t\to t+5} \;-\; \frac{1}{N_t}\sum_{j\in\mathcal S_t} r_{j,t\to t+5},
\end{equation}
the average running, with equal weights, over the stocks of date $t$ with a valid target. It sums to
zero across stocks every day; the risk-free rate cancels. All reporting is on this basis. \dec{2026-10-05}
The alternatives considered were raw or excess returns, cross-sectionally standardised or ranked returns,
volatility-scaled returns, and factor-residual (CAPM or Fama--French) returns.

\begin{whybox}
\begin{enumerate}
\item \textbf{It is the quantity the question is about.} The thesis asks whether the regime changes
  \emph{which stocks beat which}. The day's average return is the same number for every stock, carries no
  cross-sectional information, and predicting it is market timing.
\item \textbf{It keeps the gate's information and the target's apart.} The gate is fitted on the market
  series. A raw target's largest component is that same market move, so a regime that merely predicts the
  market's direction would look as if it ``helped the experts''. With the common move removed, any gain
  from the gate must come from regime-dependent \emph{cross-sectional} structure. This is the same
  attribution logic that led to freezing the gate.
\item \textbf{It makes the model's assumptions defensible.} The likelihood treats stocks as independent given
  the regime and their characteristics. For raw returns that is badly false: on a $-4\%$ day almost every
  stock falls together. Removing the common move removes the dominant source of dependence.
\item \textbf{Training and evaluation measure the same thing.} Rank IC, decile sorts and equal-weighted
  long-short returns all ignore a number added to every stock on a date, so training on the demeaned
  target spends no capacity on variation the metrics ignore.
\item \textbf{It keeps dispersion}, unlike standardising or ranking. Cross-sectional dispersion rises in stress;
  dividing by it would erase information the experts' noise scales should express.
\item \textbf{It suits a Gaussian likelihood}, unlike ranks.
\item \textbf{Equal weights in the mean} make the target match an equal-weighted, dollar-neutral long-short
  book, and keep it from being driven by a handful of mega-caps. A factor-residual target would bake one
  particular risk model into the label.
\end{enumerate}
\giveup{any claim about market timing (deliberately out of scope) and raw-return reporting.}
\end{whybox}

Note the two different ``markets'' in the pipeline: the target is demeaned by the
\emph{equal-weighted} mean of the universe; the gate and the market-model features use the
\emph{value-weighted} index (Section~\ref{sec:gateinput}). The first defines ``relative'' for an
equal-weighted book; the second measures the state of the market, which large caps drive.
